\section*{布朗运动理论}
\phantomsection
\addcontentsline{toc}{section}{布朗运动理论}

这一理论在当代科学的发展中占有特殊的地位，因为它见证了物理问题与"纯"数学之间持续而富有成果的交流。布朗运动由植物学家布朗于 1829 年发现，十九世纪中众多物理学家曾对它进行过密集的研究，\(^{3}\) 但第一个令人满意的数学模型直到 1905 年才由爱因斯坦发明。在粒子沿直线运动的简单情形中，爱因斯坦的基本假设表述如下：若 \(x(t)\) 是时刻 \(t\) 粒子的横坐标，且 \(t_0 < t_1 < \ldots < t_{n-1} < t_n\)，则相继的位移 \(x(t_i) - x(t_{i-1})\)（对 \(1 \leq i \leq n\)）是独立的高斯随机变量。此处不是详细叙述促成爱因斯坦理论的 J. 佩兰那些重要实验工作的场合：就我们的目的而言，只须记住佩兰的一个评注即可：据他说，对布朗运动轨道的观察使他不由自主地想到"数学家们的没有导数的函数"。这一评注正是点燃维纳的最初火花。

另一股完全不同的思想潮流源于气体分子运动论，它由玻尔兹曼与吉布斯在 1870 至 1900 年间发展起来。考虑由 N 个质量为 m 的分子组成、处于（绝对）温度 T 的气体，并用 \(v_{1}, \ldots, v_{N}\) 记该气体 N 个分子的速度；系统的动能等于

\[
\frac {m}{2} (\mathbf {v} _ {1} ^ {2} + \dots + \mathbf {v} _ {N} ^ {2}) = 3 N k T\tag{24.1}
\]

其中 k 是玻尔兹曼常数。按照吉布斯的想法，分子间大量的碰撞使人无法精确确定分子的速度，于是方便的做法是在由方程 (24.1) 定义的 3N 维空间的球面 S 上引入一个概率律 P。"微正则"假设在于假定 P 是球面 S 上旋转不变的质量为 1 的测度。另一方面，麦克斯韦速度定律指出，分子速度分量的概率律是方差为 2kT/m 的高斯测度。E. 博雷尔似乎是最早（1914 年）指出下述事实的人：当分子数很大时，麦克斯韦定律是吉布斯假设与球面性质的一个推论。他考虑高维欧氏空间中的一个球面 S 以及 S 上旋转不变的质量为 1 的测度 P；利用基于斯特林公式的经典逼近方法，他证明 P 在一条坐标轴上的投影近似地为高斯测度。这些结果稍后由加托与莱维加以精确化{[}201 a{]}。给定整数 \(m \geq 1\) 与数 r > 0，用 \(S_{m,r}\) 记形如 \((x_{1}, \ldots, x_{m}, 0, 0, \ldots)\) 且满足 \(x_{1}^{2} + \ldots + x_{m}^{2} = r^{2}\) 的序列所成的集；再用 \(\sigma_{m,r}\) 记 \(S_{m,r}\) 上旋转不变的质量为 1 的测度。用现代语言陈述，加托与莱维的结果如下：测度序列 \(\sigma_{m,1}\) 狭收敛于原点 \((0, 0, \ldots)\) 处的单位质量，而测度序列 \(\sigma_{m,\sqrt{m}}\) 狭收敛于形如下的测度 \(\Gamma\)

\[
d \Gamma (x _ {1}, x _ {2}, \dots) = \prod_ {n = 1} ^ {\infty} d \gamma (x _ {n})
\]

（γ 是 R 上方差为 1 的高斯测度。）

前述测度 \(\Gamma\) 起着无穷维高斯测度的作用。莱维似乎曾朦胧地希望以内在的方式在所有无穷维希尔伯特空间上定义一个高斯测度。事实上，正如莱维与维纳所证明的，测度 \(\Gamma\) 在某种意义下\(^{4}\) 在 \(l^{2}\) 的自同构作用下不变；可惜，集 \(l^{2}\)——即平方可和序列 \((x_{1}, x_{2}, \ldots, x_{n}, \ldots)\) 全体——对 \(\Gamma\) 是零测的。今天我们知道，在无穷维希尔伯特空间上我们只能满足于一个高斯前测度。\(^{5}\)

本质的进展应归功于维纳：既然无穷维希尔伯特空间上不存在合理的高斯测度，那么从高斯前测度出发，借助初等运算便可能在连续函数空间上构造一个测度 w。我们将扼要说明维纳对 w 的最初构造{[}336{]}；它直接受到加托与莱维的关系式 \(\Gamma = \lim_{m \to \infty} \sigma_{m, \sqrt{m}}\) 的影响。对每个整数 \(m \geq 1\)，用 \(H_m\) 记 T = ]0,1] 上在每个区间 \(\left[\frac{k-1}{m}, \frac{k}{m}\right]\)（对 \(k = 1, 2, \ldots, m\)）上取常值的函数所成的集，用 \(\pi_m\) 记 \(R^m\) 中半径为 1 的欧氏球面上旋转不变的质量为 1 的测度。设 \(f_m\) 是 \(H_m\) 到 \(R^m\) 上的同构，它把每个函数所取的值 \(a_k\)（该函数在区间 \(\left[\frac{k-1}{m}, \frac{k}{m}\right]\) 上取常值）对应到向量 \((a_1, a_2 - a_1, \ldots, a_m - a_{m-1})\)（维纳钟爱的"微分空间"一名即由此而来）；用 \(w_m\) 记 \(H_m\) 上的测度，即 \(\pi_m\) 在 \(f_m^{-1}\) 下的象。维纳把所求的测度 w 定义为诸测度 \(w_m\) 的极限。精确地说，用 H 记 T 上带一致收敛拓扑的正则函数所成的集（我们有 \(H_m \subset H\)，对每个整数 \(m \geq 1\)）；对每个在 H 上有界且一致连续的函数 F，极限 \(A\{F\} = \lim_{m \to \infty} \int_{H_m} F(x) dw_m(x)\) 存在；维纳通过对有关各量的涨落作细致的分析得到某些估计，并重新采用丹尼尔提出的紧性论证，证明我们正处在可以应用丹尼尔扩张定理的境地。由此推出存在一个由 \(C(T)\) 承载的测度 w，使得 \(A\{F\} = \int_{C(T)} F(x) dw(x)\)。维纳接着能够证明测度 w 对应于爱因斯坦的假设，\(^{6}\) 而他的估计使他得以给佩兰关于无导数函数的评注以精确的意义：满足 \(\frac{1}{2}\) 阶李普希茨条件的函数所成的集对 w 是可忽略的（相反，对每个满足 \(0 < a < \frac{1}{2}\) 的 a，几乎每个函数都满足 a 阶李普希茨条件）。

今天我们知道维纳测度的许多构造。例如，佩利与维纳使用随机傅里叶级数（{[}243{]}，第~IX 章）：对每个实数序列 \(\mathbf{a} = (a_{n})_{n \geq 1}\) 与每个整数 \(m \geq 0\)，在 ]0, 1] 上定义函数 \(f_{m,a}\) 如下

\[
f _ {m, \mathbf {a}} (t) = a _ {1} t + 2 \sum_ {k = 2} ^ {2 ^ {m + 1}} \frac {1}{\pi k} a _ {k - 1} \sin \pi k t;
\]

可以证明，对 \(\Gamma\)-几乎所有的序列 \(\mathbf{a}\)，函数序列 \(f_{m,\mathbf{a}}\) 趋于一个连续函数 \(f_{\mathbf{a}}\)，且 \(w\) 是 \(\Gamma\) 在（几乎处处有定义的）映射 \(\mathbf{a} \mapsto f_{\mathbf{a}}\) 下的象。后来，莱维在{[}201 b 与 c{]}中给出了另一个构造。最后，卡茨、唐斯克与埃尔德什于 1950 年前后指出，如何在维纳最初的构造中把球面测度 \(\pi_m\)（位于 \(\mathbf{R}^m\) 中）换成更一般的测度。他们的结果在维纳测度与概率论极限定理之间建立了牢固的联系；这些结果将由普罗霍罗夫加以完成和系统化{[}254{]}，对这项工作我们稍后还将回过头来讨论。

这里不是分析由维纳的发现所引起的众多而重要的概率论工作的场合；今天，布朗运动只表现为马尔可夫过程最重要的例子之一。我们只提一下卡茨把维纳测度应用于求解某些抛物型偏微分方程的工作；那是对费曼在量子场论中思想的一种改编——这是数学与物理问题之间相互影响的又一个例子。

